\documentclass[a4paper,11pt,dvipdfmx]{jsarticle}

\usepackage[top=20mm, bottom=25mm, left=20mm, right=20mm]{geometry}
\usepackage{url}
\usepackage{graphicx}
\usepackage{amssymb}
\usepackage{amsmath}
\usepackage{booktabs}
\usepackage{float}
\usepackage{multirow}

\title{今週の進捗報告\\[0.5em]\large EMS最適化における目的関数の違いが\\デマンド目標値とコストに与える影響の評価}
\author{神戸大学大学院システム情報学研究科 修士1年\\山崎博之}
\date{2026年6月12日}

\begin{document}

\maketitle

\section{評価の背景と目的}
エネルギーマネジメントシステム（EMS）における蓄電池充放電計画の最適化では，デマンド目標値（契約電力）の設定が年間の電気料金に大きく影響する．デマンド目標値の設定には，大きく分けて2つのアプローチが考えられる．

\begin{itemize}
    \item \textbf{案1：ピーク電力（デマンド目標値）の最小化}：目的関数を最大買電電力変数 $s^{\mathrm{BY}}_{\mathrm{MAX}}$ の最小化のみとし，年間を通じた買電電力のピークを可能な限り抑制する．
    \item \textbf{案2：トータルコストの最小化}：目的関数を基本料金と電力量料金の総和とし，ピーク抑制と電力量料金のバランスを同時に最適化する．
\end{itemize}

案1はデマンド目標値の最小化に特化するため，基本料金の削減には有利であるが，ピーク抑制のために蓄電池を過度に温存・使用することで電力量料金が増加する可能性がある．一方，案2はトータルコストを直接最小化するため，ピーク値をある程度許容し，電力量料金の削減との最適なバランスを探索する．

この2つのアプローチにおける数理的なトレードオフを，年間一括最適化の計算結果に基づいて定量的に評価する．

\section{評価基盤の拡張}
以前から用いている年間一括LPモデルに対し，目的関数を切り替えて評価できる枠組みを構築した．具体的には，以下の2つの目的関数アプローチを設定可能とした．

\begin{itemize}
    \item \textbf{ピーク電力（デマンド目標値）最小化（案1）}：目的関数を $s^{\mathrm{BY}}_{\mathrm{MAX}}$ の最小化のみに設定する．
          \begin{equation}
              \text{Minimize} \quad s^{\mathrm{BY}}_{\mathrm{MAX}}
          \end{equation}
    \item \textbf{トータルコスト最小化（案2）}：従来と同一の目的関数である，基本料金と電力量料金の総和を最小化する．
          \begin{equation}
              \text{Minimize} \quad w_{\mathrm{basic}} \cdot s^{\mathrm{BY}}_{\mathrm{MAX}} + \sum_{k=0}^{K-1} \left( p^{\mathrm{BY}}_{k} \cdot s^{\mathrm{BY}}_{k} \cdot \Delta t \right)
          \end{equation}
\end{itemize}
ここで，$s^{\mathrm{BY}}_{\mathrm{MAX}}$ は年間最大買電電力 [kW]（デマンド目標値），$w_{\mathrm{basic}}$ は基本料金単価に基づく年間基本料金単価 [円/kW]（$2,829.60 \times 0.85 \times 12 = 28,861.92$\,円/kW），$K$ は年間総ステップ数（$17,520$），$p^{\mathrm{BY}}_{k}$ はステップ $k$ における電力量料金単価 [円/kWh]（$21.51$\,円/kWh），$s^{\mathrm{BY}}_{k}$ はステップ $k$ における買電電力 [kW]，$\Delta t$ は時間解像度（$0.5$\,h）である．


制約条件は両モードで完全に共通であり，蓄電池パラメータ（定格容量 860.0\,kWh，最大充放電出力 400.0\,kW，片側充放電効率 0.98），需要・PVデータ（北海道十勝地方の施設における2024年の実測データ，30分粒度・17,520ステップ），および電気料金体系（固定電力量単価 21.51\,円/kWh，基本料金単価 2,829.60\,円/kW・月，力率割引 85\%）もすべて同一条件で統一されている．

なお，案1のモードにおいても，得られた最大買電電力および各ステップの買電電力から，基本料金と電力量料金を事後的に算出し，年間トータルコストを計算した．

\section{計算結果と考察}
表\ref{tab:comparison_results}に，両モードで得られた年間一括最適化の計算結果を示す．

\begin{table}[H]
    \centering
    \caption{目的関数の違いによる年間最適化結果の比較}
    \label{tab:comparison_results}
    \begin{tabular}{lrr}
        \toprule
        評価項目          & 案1（ピーク最小化）          & 案2（トータルコスト最小化）      \\
        \midrule
        最大買電電力 [kW]   & \textbf{156.05}     & 166.83              \\
        年間電力量料金 [円]   & 11,378,938          & \textbf{11,353,004} \\
        年間基本料金 [円]    & \textbf{4,503,779}  & 4,815,088           \\
        年間トータルコスト [円] & \textbf{15,882,717} & 16,168,092          \\
        \midrule
        計算時間 [秒]      & 604.33              & 600.12              \\
        理論上の最適解との最大誤差 & 0.04\%              & 0.02\%              \\
        \bottomrule
    \end{tabular}
\end{table}

\subsection*{ピーク電力と基本料金}
案1は最大買電電力を156.05\,kWに抑制し，案2の166.83\,kWに対して10.78\,kW（6.5\%）の削減を達成した．これに伴い，年間基本料金は案2の4,815,088\,円に対し4,503,779\,円となり，311,309\,円（6.5\%）の削減となった．

\subsection*{電力量料金のトレードオフ}
一方，案1の年間電力量料金は11,378,938\,円であり，案2の11,353,004\,円に対して25,934\,円（0.23\%）の増加にとどまった．電力量料金への影響は極めて限定的であった．

\subsection*{トータルコスト}
案1の年間トータルコストは15,882,717\,円，案2は16,168,092\,円であり，案1の方が285,375\,円（1.76\%）低い結果となった．これは，案1におけるピーク抑制による基本料金の削減幅（311,309\,円）が，電力量料金の増加分（25,934\,円）を大幅に上回ったためである．

\subsection*{結果の解釈}
案2のモデルは基本料金と電力量料金の加重和を同時に最小化するため，理論上は案1以上のトータルコストにはならないはずであるが，今回のシミュレーションでは案1の方がトータルコストが低い結果となった．重要な知見として，固定単価プランにおいてはピーク最小化を優先する案1が，電力量料金の増加をわずかに許容しつつも基本料金を大幅に削減し，結果的にトータルコストの観点でも案2と同等以上の性能を示すことが確認された．このことは，固定単価プランの下では，デマンド目標値の設定においてピーク抑制を主目的とするアプローチが経済性の面でも合理的であることを示唆している．

\end{document}
